\subsection{Study 2 Results}
We analyzed the replication dataset from the 0318, 0319, and 0322 collection batches. The report loaded 142 participant CSV files and retained 140 participants after excluding two attention-check failures, yielding 840 observations with a complete $3 \times 2$ within-subject grid across interaction design and detail presentation. Trust items were normalized to a 1--7 scale and averaged into a composite trust score, and the nine interaction items were averaged into an interaction experience score.

\noindent\textbf{Results: Trust.}
We first took an overview of a linear model predicting composite trust from visualization technique, interaction design, detail presentation, and their interactions. Overall, interaction design significantly affected trust, $F(2, 822)=18.73$, $p<.001$, and visualization technique also had a smaller main effect, $F(2, 822)=3.59$, $p=.028$. Detail presentation was not significant, $F(1, 822)=2.44$, $p=.118$, and none of the two-way or three-way interactions were significant (all $p>.48$). Across interaction designs, click had the highest estimated trust ($M=4.92$, $SE=0.07$), followed by hover ($M=4.85$, $SE=0.07$), while animation was substantially lower ($M=4.35$, $SE=0.07$). Trust was numerically higher for \textit{show one} ($M=4.77$, $SE=0.06$) than \textit{show all} ($M=4.64$, $SE=0.06$), but this difference was not reliable. Descriptively, trust was highest for PI ($M=4.81$), followed by Ensemble ($M=4.76$), and lowest for PI+Ensemble ($M=4.55$). At the condition level, PI with click and \textit{show one} yielded the highest mean trust ($M=5.13$, $SE=0.17$), whereas PI+Ensemble with animation and \textit{show all} yielded the lowest ($M=3.85$, $SE=0.20$).

\noindent\textbf{Takeaway.}
Allowing participants to reveal details through direct interaction increased trust relative to automatic animation, and this advantage generalized across visualization techniques.

\noindent\textbf{Results: Interaction Experience.}
We next examined the interaction experience composite with the same model structure. Interaction design again had a strong main effect, $F(2, 822)=48.08$, $p<.001$, and detail presentation also mattered, $F(1, 822)=5.32$, $p=.021$. Visualization technique was not significant, $F(2, 822)=0.46$, $p=.630$, and none of the interaction terms were significant (all $p>.84$). Across interaction designs, click again produced the highest estimated interaction experience ($M=3.98$, $SE=0.07$), followed by hover ($M=3.87$, $SE=0.07$), while animation was markedly lower ($M=3.10$, $SE=0.07$). Unlike trust, interaction experience was reliably higher for \textit{show one} ($M=3.74$, $SE=0.06$) than for \textit{show all} ($M=3.56$, $SE=0.06$). Descriptively, the highest condition mean again occurred for PI with click and \textit{show one} ($M=4.19$, $SE=0.15$), and the lowest for PI+Ensemble with animation and \textit{show all} ($M=2.85$, $SE=0.17$).

\noindent\textbf{Takeaway.}
Direct control over the transition improved the interaction experience more than automatic playback, and revealing one detail at a time further improved that experience.

\noindent\textbf{Results: Generalizability Across Visualization Techniques.}
Finally, we examined whether the effects of interaction design depended on the visualization technique. For both trust and interaction experience, the interaction between visualization technique and interaction design was not significant (trust: $F(4, 822)=0.87$, $p=.484$; interaction experience: $F(4, 822)=0.33$, $p=.859$), and no higher-order interaction involving detail presentation was significant. This indicates that the relative ordering of click and hover above animation was stable across PI, Ensemble, and PI+Ensemble. In other words, the benefit was not driven merely by seeing a transition unfold: automatic animation, which preserved the transition without requiring user action, consistently produced lower trust and lower interaction experience than the interactive conditions.

\noindent\textbf{Takeaway.}
The benefits of interaction design were general rather than technique-specific, and they appear to be caused by user control rather than by the presence of transition alone.
